\begin{workedexample}[Worked Example: T1 recovery]
Longitudinal magnetization recovers as
$M_z(t) = M_0(1 - e^{-t/T_1})$. After a $90^{\circ}$ pulse, at $t = T_1$ the
recovery is $1 - e^{-1} = 63\%$ of $M_0$. Short-$T_1$ tissue (e.g.\ fat) recovers
quickly and looks bright on $T_1$-weighted images (short TR); long-$T_1$ tissue
(e.g.\ CSF) recovers slowly and looks dark.
\end{workedexample}

\begin{keytakeaways}
\begin{itemize}[leftmargin=1.4em,itemsep=2pt]
\item $T_1$ (longitudinal recovery) and $T_2$ (transverse decay) are tissue-specific; $T_2^{*} \le T_2$.
\item $M_z(t) = M_0(1 - e^{-t/T_1})$; $M_{xy}(t) = M_0 e^{-t/T_2}$.
\item Gadolinium shortens $T_1$, brightening enhancing tissue.
\end{itemize}
\end{keytakeaways}

\begin{exercises}
\item Fraction of $T_1$ recovery at $t = T_1$?
\item Fraction of transverse signal left at $t = T_2$?
\item Does gadolinium predominantly shorten $T_1$ or $T_2$ clinically?
\end{exercises}

\furtherreading{Nishimura~\cite{nishimura}; Haacke~\cite{haacke} (ch.~4).}
